\documentclass[a4paper,11pt]{ctexart}

\usepackage[top=1.8cm,bottom=1.8cm,left=2.1cm,right=2.1cm]{geometry}
\usepackage{amsmath,amssymb}
\usepackage{enumitem}
\usepackage{fancyhdr}
\usepackage{lastpage}
\usepackage{array}

\pagestyle{fancy}
\fancyhf{}
\renewcommand{\headrulewidth}{0pt}
\fancyfoot[C]{\small 第~\thepage~页\quad 共~\pageref{LastPage}~页}

\linespread{1.28}
\setlength{\parindent}{0pt}

\newcommand{\blank}[1][2.6cm]{\underline{\hspace{#1}}}

\begin{document}

\begin{center}
  {\Large\bfseries 华中科技大学 2023～2024 学年第一学期}\\[0.35em]
  {\Large\bfseries 《高等数学（A）（上）》期中考试}\\[0.6em]
  {\normalsize
    考试方式：闭卷\quad
    考试日期：2023/11/11\quad
    考试时间：2:30--5:00 PM\\[0.35em]
    院系：\blank[2.8cm]\quad
    班级：\blank[2.8cm]\\[0.3em]
    学号：\blank[2.8cm]\quad
    姓名：\blank[2.8cm]
  }
\end{center}

\begin{center}
  \renewcommand{\arraystretch}{1.35}
  \begin{tabular}{|c|*{5}{wc{1.55cm}|}}
    \hline
    题号 & 一 & 二 & 三 & 四 & 总分 \\
    \hline
    分数 & & & & & \\
    \hline
  \end{tabular}
\end{center}

\vspace{0.7em}

\section*{第一部分\quad 期中考试试题}

\textbf{一、填空题（每题 4 分，共 28 分）}

\begin{enumerate}[label=\arabic*., leftmargin=2em, itemsep=0.75em, topsep=0.5em]
  \item 集合 $E=\left\{\dfrac{1}{m}+\dfrac{1}{n} : m,n=1,2,\cdots\right\}$，则 $\sup E=\blank$；$\inf E=\blank$.

  \item 用 $\varepsilon$ 语言写出 $\lim\limits_{n\to\infty}a_n\neq 0$ 的定义：\blank[6.5cm].

  \item 极限 $\lim\limits_{n\to\infty}\sqrt[n]{2^n+3^n}=\blank$.

  \item 已知 $\sqrt{1+\tan x}-\sqrt{1+\sin x}\sim a x^c\ (x\to 0)$，则 $a=\blank[2.0cm]$，$c=\blank[2.0cm]$.

  \item $f(x)=\begin{cases} x, & x\in\mathbb{Q}, \\ -x, & x\notin\mathbb{Q}, \end{cases}$ 指出 $f(x)$ 的间断点和间断点的类型：\blank[5.5cm].

  \item 设方程 $x^4+y^4=a^4\ (a>0\text{ 为常数})$ 在 $y\neq 0$ 时确定一个隐函数 $y(x)$，则 $y'=\blank$.

  \item 叙述海涅（Heine）定理：$\lim\limits_{x\to x_0}f(x)=A$ 的充要条件是 \blank[6.0cm].
\end{enumerate}

\vspace{0.8em}
\textbf{二、计算（每题 6 分，共 24 分）}

\begin{enumerate}[label=\arabic*., start=8, leftmargin=2em, itemsep=0.9em, topsep=0.5em]
  \item 设函数 $f(x)$ 满足 $f(7)=\dfrac{1}{17}$ 且 $f'(7)=17$，求极限 $\lim\limits_{x\to 7}\left[\dfrac{f(x)}{f(7)}\right]^{\frac{1}{\ln\frac{x}{7}}}$.

  \item 求极限 $\lim\limits_{x\to 0}\dfrac{1}{x}\bigl[(1+x)^{\frac{1}{x}}-e\bigr]$.

  \item 圆的渐开线的参数方程为 $\begin{cases} x=\cos t+t\sin t \\ y=\sin t-t\cos t \end{cases}$，其中 $t\in\left(0,\dfrac{\pi}{2}\right)$，求 $\dfrac{d^2y}{dx^2}$.

  \item 设 $f(x)=(x^2+x)e^{3x}$，求 $f^{(2023)}(x)$.
\end{enumerate}

\vspace{0.8em}
\textbf{三、证明题（每题 7 分，共 28 分）}

\begin{enumerate}[label=\arabic*., start=12, leftmargin=2em, itemsep=0.9em, topsep=0.5em]
  \item 用 $\varepsilon$ 语言证明：
  \[
    \lim_{x\to 1}\frac{x^2+1}{2x^2+x}=\frac{2}{3}.
  \]

  \item 设 $f$ 是定义在 $(0,+\infty)$ 上的函数，满足 $\forall x\in(0,+\infty),\ f(2x)=f(x)$，且 $\lim\limits_{x\to+\infty}f(x)=0$. 证明：$\forall x\in(0,+\infty),\ f(x)=0$.

  \item 设 $f$ 在区间 $I$ 上可导且导函数 $f'(x)$ 有界，证明：$f$ 在区间 $I$ 上一致连续.

  \item 设 $f(x)$ 是有界闭区间 $[a,b]$ 上的连续函数，且 $\forall x\in[a,b],\ f(x)\neq 0$. 证明：存在常数 $m>0$，满足 $\forall x\in[a,b],\ |f(x)|\ge m$.
\end{enumerate}

\vspace{0.8em}
\textbf{四、解答题（每题 10 分，共 20 分）}

\begin{enumerate}[label=\arabic*., start=16, leftmargin=2em, itemsep=0.9em, topsep=0.5em]
  \item 设 $f$ 是定义在 $\mathbb{R}$ 上的函数 $f(x)=\begin{cases} x^a\sin\dfrac{1}{x}, & x\neq 0, \\ 0, & x=0, \end{cases}$ 其中 $a>0$ 是常数，分别求 $a$ 的范围使得下面命题成立.
  \begin{enumerate}[label=(\arabic*), leftmargin=2.2em, itemsep=0.25em, topsep=0.3em]
    \item $f$ 在 $0$ 处可导并求 $f'(0)$.
    \item $f'(x)$ 在 $0$ 处连续.
    \item $f$ 在 $0$ 处二阶可导.
  \end{enumerate}

  \item 设 $\{x_n\}$ 为数列，若存在常数 $M$，满足对 $\forall n=1,2,\cdots$，都成立
  \[
    |x_2-x_1|+|x_3-x_2|+\cdots+|x_{n+1}-x_n|\le M,
  \]
  则称数列 $\{x_n\}$ 为有界变差数列.
  \begin{enumerate}[label=(\arabic*), leftmargin=2.2em, itemsep=0.25em, topsep=0.3em]
    \item 设 $x_n=\dfrac{(-1)^n}{n^2}$，验证 $\{x_n\}$ 是否为有界变差数列.
    \item 证明：有界变差数列一定是收敛数列.
    \item 举例说明收敛数列不一定是有界变差数列.
  \end{enumerate}
\end{enumerate}

\newpage
\section*{第二部分\quad 期中试题讲解与参考解答}

\textbf{一、填空题}

\begin{enumerate}[label=\arabic*., leftmargin=2em, itemsep=0.75em, topsep=0.4em]
  \item $\sup E=2$；$\inf E=0$.

  \textbf{解：}$m,n$ 都取 $1$ 时，$E$ 取到最大值，所以 $\sup E=2$；又 $\forall m,n\ge 1,\ \dfrac{1}{m}+\dfrac{1}{n}>0$ 且 $\lim\limits_{k\to\infty}\left(\dfrac{1}{k}+\dfrac{1}{k}\right)=0$，故 $\inf E=0$.

  \item $\exists\,\varepsilon_0>0,\ \forall N\in\mathbb{N}^+,\ \exists\,n_0>N,\ |a_{n_0}|\ge\varepsilon_0$.

  \item $3$.

  \textbf{解：}利用夹挤原理：
  \[
    3=\sqrt[n]{3^n}<\sqrt[n]{2^n+3^n}<\sqrt[n]{2\cdot 3^n}=3\sqrt[n]{2},
  \]
  $\lim\limits_{n\to\infty}3=3$，$\lim\limits_{n\to\infty}3\sqrt[n]{2}=3$，所以 $\lim\limits_{n\to\infty}\sqrt[n]{2^n+3^n}=3$.

  \item $a=\dfrac{1}{4}$，$c=3$（即 $ax^c=\dfrac{x^3}{4}$）.

  \textbf{解：}因为 $\sqrt{1+\tan x}-\sqrt{1+\sin x}\sim ax^c\ (x\to 0)$，故
  \begin{align*}
    \sqrt{1+\tan x}-\sqrt{1+\sin x}
    &= \frac{\tan x-\sin x}{\sqrt{1+\tan x}+\sqrt{1+\sin x}} \\
    &\sim \frac{x+\dfrac{x^3}{3}-\left(x-\dfrac{x^3}{6}\right)+o(x^3)}{2}
    = \frac{x^3}{4}+o(x^3)\quad (x\to 0),
  \end{align*}
  所以 $ax^c=\dfrac{x^3}{4}$，即 $a=\dfrac{1}{4},\ c=3$.

  \item 间断点为一切 $x\neq 0$，且均为第二类间断点.

  \textbf{解：}取任意 $\{x_n\}$ 满足 $\lim\limits_{n\to\infty}x_n=0$，有 $|f(x_n)|=|x_n|\to 0$，故 $\lim\limits_{n\to\infty}f(x_n)=0=f(0)$，因此由 Heine 定理可知，$0$ 为 $f(x)$ 的连续点；

  $\forall x_0\neq 0$，取有理点列 $\{y_n\}$ 满足 $y_n\in U^\circ(x_0)\cap\mathbb{Q},\ \lim\limits_{n\to\infty}y_n=x_0$，则 $\lim\limits_{n\to\infty}f(y_n)=\lim\limits_{n\to\infty}y_n=x_0$；

  取无理点列 $\{z_n\}$ 满足 $z_n\in U^\circ(x_0)\setminus\mathbb{Q},\ \lim\limits_{n\to\infty}z_n=x_0$，则 $\lim\limits_{n\to\infty}f(z_n)=-\lim\limits_{n\to\infty}z_n=-x_0\neq x_0$，

  所以 $\lim\limits_{x\to x_0}f(x)$ 不存在，即一切 $x_0\neq 0$ 均为 $f(x)$ 的第二类间断点.

  \item $-\dfrac{x^3}{y^3}$.

  \textbf{解：}在方程两边对 $x$ 求导可得：$4x^3+4y^3y'=0$，解得 $y'=-\dfrac{x^3}{y^3}$.

  \item $\forall \{x_n\}\subset U^\circ(x_0)$（或函数定义域内且 $x_n\neq x_0$），若 $\lim\limits_{n\to\infty}x_n=x_0$，则 $\lim\limits_{n\to\infty}f(x_n)=A$.
\end{enumerate}

\vspace{0.6em}
\textbf{二、计算}

\begin{enumerate}[label=\arabic*., start=8, leftmargin=2em, itemsep=0.9em, topsep=0.4em]
  \item \textbf{解：}
  \[
    \lim_{x\to 7}\left[\frac{f(x)}{f(7)}\right]^{\frac{1}{\ln\frac{x}{7}}}
    = e^{\lim\limits_{x\to 7}\frac{\ln f(x)-\ln f(7)}{\ln x-\ln 7}}
    = e^{\lim\limits_{x\to 7}\frac{\ln f(x)-\ln f(7)}{x-7}\big/\lim\limits_{x\to 7}\frac{\ln x-\ln 7}{x-7}}
    = e^{\frac{f'(7)}{f(7)}\big/\frac{1}{7}}
    = e^{2023}.
  \]

  \item \textbf{解：}
  \begin{align*}
    \lim_{x\to 0}\frac{1}{x}\bigl[(x+1)^{\frac{1}{x}}-e\bigr]
    &= e\lim_{x\to 0}\frac{e^{\frac{\ln(x+1)}{x}-1}-1}{x}
    = e\lim_{x\to 0}\frac{\dfrac{\ln(x+1)}{x}-1}{x} \\
    &= e\lim_{x\to 0}\frac{\ln(x+1)-x}{x^2}
    = e\lim_{x\to 0}\frac{x-\dfrac{x^2}{2}+o(x^2)-x}{x^2}
    = -\frac{e}{2}.
  \end{align*}

  \item \textbf{解：}
  \[
    \frac{dy}{dx}=\frac{\dfrac{dy}{dt}}{\dfrac{dx}{dt}}=\frac{\cos t-(\cos t-t\sin t)}{-\sin t+\sin t+t\cos t}=\frac{t\sin t}{t\cos t}=\tan t,
  \]
  \[
    \frac{d^2y}{dx^2}=\frac{\dfrac{d(\tan t)}{dt}}{\dfrac{dx}{dt}}=\frac{\sec^2 t}{t\cos t}=\frac{1}{t\cos^3 t}.
  \]

  \item \textbf{解：}利用 Leibniz 求导公式：
  \begin{align*}
    f^{(2023)}(x)
    &= \sum_{k=0}^{2023} C_{2023}^k (x^2+x)^{(k)}(e^{3x})^{(2023-k)} \\
    &= C_{2023}^0 (x^2+x)(e^{3x})^{(2023)} + C_{2023}^1 (x^2+x)'(e^{3x})^{(2022)} + C_{2023}^2 (x^2+x)''(e^{3x})^{(2021)} \\
    &= 3^{2023}(x^2+x)e^{3x} + 2023\cdot 3^{2022}(2x+1)e^{3x} + 2\cdot 3^{2021} C_{2023}^2 e^{3x} \\
    &= 3^{2021}e^{3x}\bigl[9x^2+(2023\cdot 6+9)x+2023\cdot 2025\bigr].
  \end{align*}
\end{enumerate}

\vspace{0.6em}
\textbf{三、证明题}

\begin{enumerate}[label=\arabic*., start=12, leftmargin=2em, itemsep=0.9em, topsep=0.4em]
  \item \textbf{证：}$\forall\varepsilon>0$，取 $\delta=\min\left\{\dfrac{2\varepsilon}{3},\,\dfrac{1}{2}\right\}$，当 $0<|x-1|<\delta$ 时（有 $\dfrac{1}{2}<x<\dfrac{3}{2}$），
  \[
    \left|\frac{x^2+1}{2x^2+x}-\frac{2}{3}\right|
    =\frac{|x^2+2x-3|}{3|2x^2+x|}
    =\frac{|x-1|\cdot|x+3|}{3|x|\cdot|2x+1|}
    <\frac{|x-1|\cdot\dfrac{9}{2}}{3\cdot\dfrac{1}{2}\cdot 2}
    =\frac{3}{2}|x-1|<\varepsilon.
  \]
  所以 $\lim\limits_{x\to 1}\dfrac{x^2+1}{2x^2+x}=\dfrac{2}{3}$.

  \item \textbf{证明：}$\forall x\in(0,+\infty)$，$f(x)=0$.

  \textbf{证法 1：}$\forall x\in(0,+\infty)$，记 $a_n=f(2^n x)$. 由 $f(2x)=f(x)$ 知 $\{a_n\}$ 为常值数列；由 $\lim\limits_{x\to+\infty}f(x)=0$ 知 $\lim\limits_{n\to\infty}a_n=\lim\limits_{n\to\infty}f(2^n x)=0$，所以 $\forall n,\ a_n=0$，故 $a_0=f(x)=0$.

  \textbf{证法 2（反证法）：}若 $\exists\,x_0\in(0,+\infty)$，s.t\ $f(x_0)\neq 0$，则 $f(x_0)=f(2x_0)=\cdots=f(2^n x_0)\neq 0$，$\{f(2^n x_0)\}$ 为非零常值数列；由 $\lim\limits_{x\to+\infty}f(x)=0$ 知 $\lim\limits_{n\to\infty}f(2^n x_0)=0$，这与 $f(2^n x_0)=f(x_0)\neq 0$ 矛盾，所以 $\forall x\in(0,+\infty),\ f(x)=0$.

  \item \textbf{证：}$f'(x)$ 在区间 $I$ 上有界，即 $\exists M>0,\ \forall x\in I,\ |f'(x)|\le M$，则 $\forall\varepsilon>0$，取 $\delta=\dfrac{\varepsilon}{M}$，$\forall x,y\in I$ 且 $|x-y|<\delta$ 时（若 $x=y$ 则显然成立；若 $x\neq y$），由 Lagrange 中值定理存在 $\xi$ 介于 $x,y$ 之间，使得
  \[
    |f(x)-f(y)|=|f'(\xi)(x-y)|\le M|x-y|<M\delta=\varepsilon,
  \]
  因此由一致连续的定义可知 $f(x)$ 在 $I$ 上一致连续.

  \item \textbf{证：}$\forall x\in[a,b],\ f(x)\neq 0$，由连续函数的介值定理知 $f(x)$ 在区间 $[a,b]$ 上不变号，故不妨设 $\forall x\in[a,b],\ f(x)>0$（若 $f(x)<0$，可考虑 $-f(x)$）. 由最值定理知 $f(x)$ 在区间 $[a,b]$ 上可取到最小值 $f(x_0)$ 且 $f(x_0)>0$，取 $m=f(x_0)$ 即可.
\end{enumerate}

\vspace{0.6em}
\textbf{四、解答题}

\begin{enumerate}[label=\arabic*., start=16, leftmargin=2em, itemsep=0.9em, topsep=0.4em]
  \item \textbf{解：}(1) $f'(0)=\lim\limits_{x\to 0}\dfrac{x^a\sin\frac{1}{x}-0}{x-0}=\lim\limits_{x\to 0}x^{a-1}\sin\dfrac{1}{x}$ 存在，此时 $a>1$，且 $f'(0)=0$；

  (2) $x\neq 0$ 时，$f'(x)=a x^{a-1}\sin\dfrac{1}{x}-x^{a-2}\cos\dfrac{1}{x}$. $f'(x)$ 在 $0$ 点连续，当且仅当 $\lim\limits_{x\to 0}f'(x)=\lim\limits_{x\to 0}\left(a x^{a-1}\sin\dfrac{1}{x}-x^{a-2}\cos\dfrac{1}{x}\right)=f'(0)$，此时 $a>2$；

  (3) 当 $a>1$ 时 $f$ 在 $\mathbb{R}$ 上可导，故
  \[
    f''(0)=\lim_{x\to 0}\frac{a x^{a-1}\sin\dfrac{1}{x}-x^{a-2}\cos\dfrac{1}{x}-0}{x-0}
    =\lim_{x\to 0}\left(a x^{a-2}\sin\frac{1}{x}-x^{a-3}\cos\frac{1}{x}\right)
  \]
  存在，此时 $a>3$（且 $f''(0)=0$）.

  \item \textbf{解：}(1)
  \begin{align*}
    |x_2-x_1|+|x_3-x_2|+\cdots+|x_{n+1}-x_n|
    &= \left(1+\frac{1}{2^2}\right)+\left(\frac{1}{2^2}+\frac{1}{3^2}\right)+\cdots+\left(\frac{1}{n^2}+\frac{1}{(n+1)^2}\right) \\
    &\le 2\left(1+\frac{1}{2^2}+\frac{1}{3^2}+\cdots+\frac{1}{n^2}+\frac{1}{(n+1)^2}\right) \\
    &< 2\left(1+\frac{1}{1\cdot 2}+\frac{1}{2\cdot 3}+\cdots+\frac{1}{n(n+1)}\right) < 4,
  \end{align*}
  即 $\{x_n\}$ 为有界变差数列；

  (2) 记 $a_n=|x_2-x_1|+|x_3-x_2|+\cdots+|x_{n+1}-x_n|$，则 $\{a_n\}$ 为单调非减有界数列，所以 $\{a_n\}$ 收敛. 由 Cauchy 收敛准则，$\forall\varepsilon>0,\ \exists N,\ \forall n>N,\ \forall p\in\mathbb{N}^+,\ |a_{n+p}-a_n|<\varepsilon$，又
  \[
    |x_{n+1+p}-x_{n+1}|\le |x_{n+1+p}-x_{n+p}|+|x_{n+p}-x_{n+p-1}|+\cdots+|x_{n+2}-x_{n+1}|=|a_{n+p}-a_n|<\varepsilon,
  \]
  所以 $\{x_n\}$ 为 Cauchy 列，因此收敛；

  (3) 取 $x_n=\dfrac{(-1)^n}{n}$，$\{x_n\}$ 是收敛数列，但
  \[
    |x_2-x_1|+|x_3-x_2|+\cdots+|x_{n+1}-x_n|
    =1+\frac{1}{2}+\frac{1}{2}+\frac{1}{3}+\cdots+\frac{1}{n}+\frac{1}{n+1}\to+\infty\quad (n\to\infty),
  \]
  故 $\{x_n\}$ 不是有界变差数列.
\end{enumerate}

\end{document}
